\section{Related Work}

\paragraph{Consistency Evaluation on Transformed Inputs.}
ML models should exhibit invariance to small input transformations~\citepia{42503,10.1145/3052973.3053009,8424625}.
However, this is often violated when models overfit to their training data.
For example, past work have found that translation models~\citep{belinkov2018synthetic}, NLI models~\citep{arakelyan-etal-2024-semantic}, QA/NER/sentiment models~\citep{Rychalska2019models}, etc., degrade in performance under meaning-preserving input changes.
General-purpose LMs have likewise been shown to be sensitive to minor input 
\iflongversion
transformations~\citepia{lu-etal-2022-fantastically,gonen-etal-2023-demystifying,sclar2024quantifying} or larger changes that robust models should have invariant predictions~\citepia{wu-etal-2024-reasoning,doi:10.1073/pnas.2322420121}.
Various benchmarks have likewise been developed to test model robustness~\citepia{chao2024jailbreakbench,ye-etal-2024-rotbench,jung-etal-2025-flex}.
To our knowledge, our
\else
transformations~\citepia{lu-etal-2022-fantastically,gonen-etal-2023-demystifying,sclar2024quantifying}.
Our
\fi
work is the first to show this for RMs, which is particularly significant (\S\ref{sec:background}).

\paragraph{Improving Model Robustness.}
Due to the importance of model robustness, much work has explicitly trained models to be less brittle. %
Many models are trained to be consistent with respect to data augmentation~\citepia{DBLP:journals/corr/GuR14,43405,zhang2019theoretically,9156853,tack2022consistency}.
Past work also trained LMs to be robust on various tasks~\citepia{zheng-etal-2021-consistency,zhou-etal-2022-prompt,yan-etal-2024-contrastive,zhou-etal-2024-fairer}.
Our work inherits these ideas to train more robust RMs.
Similar to us, \citet{shen2024the} also trained regularized RMs, though with different objectives.
